\section*{अध्याय \ref{ch:b2:magma-emplacement} --- मैग्मा और उनका स्थापन}

\begin{solution}{exo:b2:magma-emplacement:1}
$T = \qty{1073}{K}$। $w = 1$ के लिए, $\ln w = 0$: $\log_{10}\eta = -3.545 +
9601/(1073 - 195.7) = -3.545 + 10.94 = 7.40$। $w = 4$ के लिए, $\ln 4 = 1.386$:
$A = -3.545 + 1.155 = -2.390$, $B = 9601 - 3283 = 6318$, $C = 195.7 + 44.7 =
240.4$, और $\log_{10}\eta = -2.390 + 6318/832.6 = -2.390 + 7.588 = 5.20$।
श्यानता इस गुणक से घटती है: $10^{7.40 - 5.20} = 10^{2.20} \approx 160$।
\end{solution}

\begin{solution}{exo:b2:magma-emplacement:2}
रायोलाइट: $p_s = (6/0.44)^2 = \qty{186}{MPa}$, जो इस गहराई पर है: $\num{186e6}/(2600\times9.81)
= \qty{7.3}{km}$। बेसाल्ट: $p_s = (0.5/0.34)^2 = \qty{2.16}{MPa}$, जो
$\num{2.16e6}/(2600\times9.81) = \qty{85}{m}$ पर है। रायोलाइट भूपर्पटी में गहराई पर ही बुलबुले
बनाने लगता है, बेसाल्ट केवल पृष्ठ के निकट (यदि उसमें CO$_2$ न हो)।
\end{solution}

\begin{solution}{exo:b2:magma-emplacement:3}
$(1150 + 50)/2 = \qty{600}{\celsius}$ और $(1150 + 400)/2 = \qty{775}{\celsius}$। ठंडी
चट्टान वाला \omterm{def:b2:magma-emplacement:bodies}{डाइक} अपने संपर्क पर अधिक तीव्रता से शीतित होता है, और उसी का
\omterm{def:b2:magma-emplacement:aureole}{शीतित किनारा} अधिक चौड़ा होता है।
\end{solution}

\begin{solution}{exo:b2:magma-emplacement:4}
(क) एक \omterm{def:b2:magma-emplacement:bodies}{सिल}, संस्तरण के समांतर। (ख) एक \omterm{def:b2:magma-emplacement:bodies}{डाइक}, संस्तरण को काटता हुआ।
(ग) एक \omterm{def:b2:magma-emplacement:bodies}{बैथोलिथ}, \qty{100}{km^2} से बड़ा और अनेक अंतर्वेधनों से एकत्रित। (घ) एक
\omterm{def:b2:magma-emplacement:bodies}{लैकोलिथ}, जिसने अपनी छत को गुंबदाकार उठाया है। (ङ) एक \omterm{def:b2:magma-emplacement:bodies}{डाइक}: ग्रेनाइट में कोई परतें नहीं होतीं,
इसलिए उसमें कोई भी चादर \omterm{def:b2:magma-emplacement:bodies}{डाइक} है, उसका अभिविन्यास कुछ भी हो।
\end{solution}

\begin{solution}{exo:b2:magma-emplacement:5}
$(1 - 1.35\phi)^{-2.5}$: $\phi = 0.05$ के लिए, $0.9325^{-2.5} = 1.19$, इसलिए
\qty{119}{Pa.s}, जबकि $1 + 2.5\times0.05 = 1.125$, यानी \qty{113}{Pa.s}, आइंस्टाइन के
सूत्र से (निकट, जैसा तनु सीमा में होना चाहिए)। $0.20$ के लिए: $0.73^{-2.5} =
2.20$, \qty{220}{Pa.s}; $0.40$ के लिए: $0.46^{-2.5} = 6.97$, \qty{697}{Pa.s}; $0.60$ के लिए:
$0.19^{-2.5} = 63.6$, लगभग \qty{6360}{Pa.s}। जाम $\phi = 1/1.35 = 0.74$ पर।
\end{solution}

\begin{solution}{exo:b2:magma-emplacement:6}
$G = 300\times9.81 = \qty{2943}{Pa/m}$। (क) $\bar u = 2943\times1^2/(12\times100) =
\qty{2.45}{m/s}$; $q = \qty{2.45}{m^2/s}$; $Re = 2700\times2.45\times1/100 = 66$, स्तरीय;
$t = 10000/2.45 = \qty{4080}{s}$, लगभग \qty{1.1}{h}। (ख) $\bar u =
2943/(12\times\num{e6}) = \qty{2.45e-4}{m/s}$, $t = \qty{4.08e7}{s} = 472$ दिन;
$w = \qty{10}{m}$ के साथ $\bar u$ सौ गुना बड़ा है, \qty{0.0245}{m/s}, और $t =
\qty{4.08e5}{s} = 4.7$ दिन। (ग) \qty{1}{m} का \omterm{def:b2:magma-emplacement:bodies}{डाइक} लगभग
$0.5^2/\num{e-6} = \qty{2.5e5}{s}$ में, 2.9 दिन में, ठंडा होता है; \qty{10}{m} का \omterm{def:b2:magma-emplacement:bodies}{डाइक} $5^2/\num{e-6} =
\qty{2.5e7}{s}$ में, 289 दिन में। बेसाल्ट (\qty{1.1}{h}) और चौड़े \omterm{def:b2:magma-emplacement:bodies}{डाइक} में
रायोलाइट (4.7 दिन) अपने \omterm{def:b2:magma-emplacement:bodies}{डाइकों} के जमने से पहले पहुँच जाते हैं;
\qty{1}{m} के \omterm{def:b2:magma-emplacement:bodies}{डाइक} में रायोलाइट (472 दिन) नहीं पहुँच सकता।
\end{solution}

\begin{solution}{exo:b2:magma-emplacement:7}
प्रति किलोग्राम \omterm{def:b2:magma-emplacement:magma}{मैग्मा}, जल के $n = 10/18.015 = \qty{0.555}{mol}$ हैं और $nRT =
0.555\times8.314\times1173 = \qty{5414}{J}$; शेष गलित घेरता है $0.99/2400 =
\qty{4.13e-4}{m^3}$। \qty{100}{MPa} पर: $V = \qty{5.41e-5}{m^3}$, अंश
$5.41/(5.41 + 41.3) = 0.12$। \qty{10}{MPa} पर: $V = \qty{5.41e-4}{m^3}$, अंश
$0.57$। \qty{0.1}{MPa} पर: $V = \qty{0.0541}{m^3}$, अंश $0.992$।
\qty{100}{MPa} वाला मान सबसे कम विश्वसनीय है: इतने दाब पर जलवाष्प सघन और
आदर्श से बहुत दूर होती है (और वास्तव में अधिकांश जल तब भी घुला रहेगा)।
\end{solution}

\begin{solution}{exo:b2:magma-emplacement:8}
(क) $2300 + 60z = 2650$ से $z = \qty{5.8}{km}$ मिलता है। (ख) गलित का घनत्व
\qty{2580}{kg/m^3} है: $z = 280/60 = \qty{4.7}{km}$, लगभग \qty{1.2}{km} ऊपर। (ग)
$z = \qty{7.8}{km}$ पर चट्टान का घनत्व $2300 + 60\times7.8 = \qty{2770}{kg/m^3}$ है; बल
$(2770 - 2650)\times9.81 = \qty{1.2}{kN}$ प्रति घन मीटर है, ऊपर की ओर।
\end{solution}

\begin{solution}{exo:b2:magma-emplacement:9}
$v = 2\times300\times9.81\times(1000)^2/(9\times\num{e19}) = \qty{6.5e-11}{m/s}$, यानी
$\num{6.5e-11}\times\num{3.156e7} = \qty{2.1e-3}{m}$ प्रति वर्ष, \qty{2.1}{mm/yr}।
\qty{10}{km} उठने में $\num{e4}/\num{2.1e-3} = \qty{4.8}{Myr}$ लगते हैं। शीतलन का समय-पैमाना
$(1000)^2/\num{e-6} = \qty{e12}{s} = \qty{32}{kyr}$ है: पिंड लगभग
$32000\times0.0021 \approx \qty{70}{m}$ उठने के बाद ठोस हो जाता है। मैग्मीय \omterm{def:b2:magma-emplacement:ascent}{डायापिर}
ऐसी भूपर्पटी को पार नहीं कर सकता; उसे अधिक गरम, अधिक नरम परिवेश चाहिए।
\end{solution}

\begin{solution}{exo:b2:magma-emplacement:10}
(क) $x_{\mathrm{OH}} = \sqrt{K_2x_m} = \sqrt{K_1K_2f}$, इसलिए $x = K_1f + \frac12\sqrt{K_1K_2f}$।
पहले पद का दूसरे से अनुपात $2\sqrt{K_1f/K_2}$ है: यह शून्य की ओर प्रवृत्त होता है
जब $f \to 0$ (वर्गमूल हावी रहता है) और $f$ के साथ असीमित बढ़ता है
(रैखिक पद हावी रहता है)। (ख) $K_1 f = \frac12\sqrt{K_1K_2 f}$ से $K_1^2f^2 =
K_1K_2f/4$ मिलता है, इसलिए $f = K_2/(4K_1)$। (ग) $f = 0.2/(\num{4e-4}) = \qty{500}{MPa}$।
\qty{100}{MPa} पर: $x_m = 0.0100$, $\frac12\sqrt{\num{e-4}\times0.2\times100} =
\frac12\times0.0447 = 0.0224$, इसलिए $x = 0.0324$ और हाइड्रॉक्सिल का हिस्सा
$0.0224/0.0324 = \qty{69}{\percent}$ है।
\end{solution}

\begin{solution}{exo:b2:magma-emplacement:11}
(क) $1000/\num{4e6} = \qty{2.5e-4}{km^3}$ प्रति वर्ष, अर्थात $\qty{2.5e5}{m^3}$ प्रति वर्ष,
या $\num{2.5e5}/\num{3.156e7} = \qty{0.0079}{m^3/s}$, लगभग आठ लीटर प्रति सेकंड।
(ख) \omterm{def:b2:magma-emplacement:bodies}{डाइक} $1000\times2\times0.8 = \qty{1600}{m^3/s}$ ले जाता है, और
$\qty{e12}{m^3}$ पहुँचाता है \qty{6.25e8}{s} में, लगभग 20 वर्षों में: यह $\num{5e-6}$ है
\qty{4}{Myr} का। (ग) हर चादर में $20\times\num{e4}\times\num{e4} = \qty{2e9}{m^3} =
\qty{2}{km^3}$ है, इसलिए 500 चादरें, औसतन हर $\num{4e6}/500 = 8000$ वर्ष में एक।
(घ) $t = (10)^2/(4\times0.65^2\times\num{0.7e-6}) = \qty{8.5e7}{s}$, 2.7 वर्ष। हर
चादर अगली के आने से हज़ारों वर्ष पहले ठोस हो जाती है: जब तक आपूर्ति
तीव्र झोंकों में न आए, \omterm{def:b2:magma-emplacement:bodies}{प्लूटॉन} कभी \omterm{def:b2:magma-emplacement:magma}{मैग्मा} का एक अकेला बड़ा पिंड नहीं होता, जैसा
बड़े \omterm{def:b2:magma-emplacement:bodies}{प्लूटॉनों} का काल-निर्धारण सुझाता है।
\end{solution}

\begin{solution}{exo:b2:magma-emplacement:12}
(क) स्थिर $x$ पर, $\ln(T - T_0) = \text{अचर} - \frac12\ln t - x^2/(4\kappa t)$, जिसका
अवकलज $-1/(2t) + x^2/(4\kappa t^2)$ शून्य होता है $t^* = x^2/(2\kappa)$ पर। तब
$e^{-x^2/(4\kappa t^*)} = e^{-1/2}$ और $\sqrt{4\pi\kappa t^*} = x\sqrt{2\pi}$, इसलिए $\Delta
T_{\max} = (T_m - T_0)\,b\,e^{-1/2}/(x\sqrt{2\pi}) = (T_m - T_0)\,b/(x\sqrt{2\pi e})$। (ख)
$\sqrt{2\pi e} = 4.13$: $x = 1100\times2/(100\times4.13) = \qty{5.3}{m}$, जो पहुँचता है
$t^* = 5.3^2/(2\times\num{0.7e-6}) = \qty{2.0e7}{s}$ के बाद, लगभग 230 दिन (केवल 2.7
डाइक-चौड़ाइयों पर समतल-स्रोत सन्निकटन स्थूल है)। (ग) तापन $b/x$ की तरह
घटता है: पतला \omterm{def:b2:magma-emplacement:bodies}{डाइक} कम ऊष्मा ले जाता है, और चट्टान केवल कुछ डाइक-चौड़ाइयों के भीतर
उल्लेखनीय रूप से गरम होती है।
\end{solution}

\begin{solution}{pb:b2:magma-emplacement:1}
\textbf{1.} $\log_{10}\eta = -4.55 + 5510/(1423 - 611) = -4.55 + 6.786 = 2.24$;
$\eta = \qty{172}{Pa.s}$।

\textbf{2.} $-4.55 + 5510/762 = -4.55 + 7.231 = 2.68$, $\eta = \qty{480}{Pa.s}$: \qty{50}{K} के लिए
गुणक $480/172 = 2.8$ है।

\textbf{3.} $172\times(1 - 1.35\times0.20)^{-2.5} = 172\times0.73^{-2.5} = 172\times2.20 =
\qty{378}{Pa.s}$।

\textbf{4.} $\phi = 1/1.35 = 0.74$। क्षेत्रीय साक्ष्य समग्र प्रवाह का अंत
\qtyrange{30}{50}{\percent} गलित-अंश पर रखते हैं, अर्थात
\qtyrange{50}{70}{\percent} क्रिस्टलों पर: पहले, क्योंकि वास्तविक क्रिस्टल एक आकार के
गोले नहीं होते।

\textbf{5.} आधार: $2800 - 2710 = \qty{90}{kg/m^3}$, \omterm{def:b2:magma-emplacement:magma}{मैग्मा} हल्का है।
अवसाद: $2450 - 2710 = \qty{-260}{kg/m^3}$, \omterm{def:b2:magma-emplacement:magma}{मैग्मा} सघन है।

\textbf{6.} $p = 2450\times9.81\times3000 = \qty{72.1}{MPa}$।

\textbf{7.} $90\times9.81 = \qty{883}{Pa/m}$।

\textbf{8.} $\bar u = 883\times2.0^2/(12\times378) = 3532/4536 = \qty{0.78}{m/s}$।

\textbf{9.} $q = 0.78\times2.0 = \qty{1.56}{m^2/s}$; \qty{1.0}{km} पर, लगभग
\qty{1560}{m^3/s}।

\textbf{10.} $Re = 2710\times0.78\times2.0/378 = 11$: स्तरीय।

\textbf{11.} $12000/0.78 = \qty{15400}{s}$, लगभग \qty{4.3}{h}।

\textbf{12.} $(1.0)^2/\num{0.7e-6} = \qty{1.43e6}{s}$, 16.5 दिन, \qty{4.3}{h} से
कहीं अधिक: \omterm{def:b2:magma-emplacement:magma}{मैग्मा} के बहते समय \omterm{def:b2:magma-emplacement:bodies}{डाइक} खुला रहता है।

\textbf{13.} \qty{3.0}{km} पर \omterm{def:b2:magma-emplacement:magma}{मैग्मा} अपने से हल्के अवसादों से मिलता है: अवसादों का
आधार उसका \omterm{def:b2:magma-emplacement:ascent}{उदासीन उत्प्लावकता का स्तर} है। उसके ऊपर \omterm{def:b2:magma-emplacement:magma}{मैग्मा} अपने परिवेश से सघन
होता, इसलिए वह संपर्क के अनुदिश किनारों की ओर एक \omterm{def:b2:magma-emplacement:bodies}{सिल} के रूप में
फैलता है।

\textbf{14.} $p_s = (0.40/0.34)^2 = \qty{1.38}{MPa}$, जो प्राप्त होता है
पृष्ठ से $\num{1.38e6}/(2450\times9.81) = \qty{58}{m}$ नीचे।

\textbf{15.} $0.34\times\sqrt{72.1} = 0.34\times8.49 = \qty{2.9}{\percent}$, उपस्थित
\qty{0.40}{\percent} से कहीं अधिक: \omterm{def:b2:magma-emplacement:bodies}{सिल} का \omterm{def:b2:magma-emplacement:magma}{मैग्मा} कोई जल निर्विलयित नहीं करता।

\textbf{16.} CO$_2$ जल से बहुत कम विलेय है: वह आरोहण में गहराई पर ही संतृप्त होकर
बुलबुले बनाती है, इसलिए बेसाल्ट की पहली गैस अधिकतर CO$_2$ होती है।

\textbf{17.} प्रति किलोग्राम, $n = 4.0/18.015 = \qty{0.222}{mol}$, $V = 0.222\times8.314\times1423/
\num{e5} = \qty{0.0263}{m^3}$; गलित घेरता है $0.996/2710 = \qty{3.68e-4}{m^3}$;
अंश $0.0263/(0.0263 + 0.000368) = 0.986$ है, \qty{98.6}{\percent}।

\textbf{18.} $(1150 + 100)/2 = \qty{625}{\celsius}$।

\textbf{19.} 10 वर्ष बाद, $2\sqrt{\kappa t} = 2\sqrt{\num{0.7e-6}\times\num{3.156e8}} =
\qty{29.7}{m}$। बिंदु मध्य से $x = \qty{60}{m}$ पर है: $T = 100 +
525\,[\operatorname{erf}(-10/29.7) + \operatorname{erf}(110/29.7)] = 100 + 525\,(-0.365 +
1.000) = \qty{433}{\celsius}$।

\textbf{20.} मध्य में, $T = 100 + 1050\operatorname{erf}(50/2\sqrt{\kappa t})$:
10 वर्ष बाद $1050\operatorname{erf}(1.68) = 1050\times0.983$, $T = \qty{1132}{\celsius}$;
50 वर्ष बाद, $2\sqrt{\kappa t} = \qty{66.5}{m}$, $1050\operatorname{erf}(0.752) = 1050\times0.713$,
$T = \qty{848}{\celsius}$।

\textbf{21.} $x = \qty{130}{m}$ पर: 300 वर्ष बाद \qty{294}{\celsius},
363 वर्ष बाद \qty{295}{\celsius} (अधिकतम), 420 वर्ष बाद \qty{295}{\celsius}।
गुप्त ऊष्मा के बिना, \omterm{def:b2:magma-emplacement:bodies}{सिल} से \qty{80}{m} ऊपर की चट्टान कभी पूरी तरह
\qty{300}{\celsius} तक नहीं पहुँचती।

\textbf{22.} $\num{9e8}\times1.772/(\num{2.5e6}\times1050) = 0.608$।

\textbf{23.} $\lambda = 0.64$: $e^{-0.4096}/(0.64\times1.6346) = 0.6639/1.0461 = 0.635$;
$0.65$: $0.6554/1.0673 = 0.614$; $0.66$: $0.6469/1.0886 = 0.594$। मान $0.608$
$0.65$ और $0.66$ के बीच है, $0.65$ के निकट: $\lambda = 0.65$।

\textbf{24.} $100 + 1050/(1 + 0.642) = 100 + 640 = \qty{740}{\celsius}$। ठोसीकरण-अग्र पर
मुक्त गुप्त ऊष्मा \omterm{def:b2:magma-emplacement:magma}{मैग्मा} की ओर को अधिक देर तक
\qty{1150}{\celsius} पर रखती है, इसलिए संपर्क को अधिक ऊष्मा मिलती है।

\textbf{25.} $(50)^2/\num{0.7e-6} = \qty{3.57e9}{s}$, 113 वर्ष।

\textbf{26.} $t = (100)^2/(16\times0.65^2\times\num{0.7e-6}) = \num{e4}/\num{4.73e-6} =
\qty{2.11e9}{s}$, इसलिए $\num{2.11e9}/\num{3.156e7}$: \omterm{def:b2:magma-emplacement:bodies}{सिल} ठोस हो जाती है
\textbf{\qty[text-series-to-math=true]{67}{} वर्षों} के बाद। यह छोटा है
$(b/2)^2/\kappa$ से, गुणक $1/(4\lambda^2) = 0.59$ से: वह समय-पैमाना केवल
परिमाण की कोटि तय करता है, और \omterm{def:b2:magma-emplacement:bodies}{सिल} की दोनों दीवारें एक साथ ठंडे आतिथेय में ऊष्मा
खींचती हैं, इतनी तेज़ी से कि गुप्त ऊष्मा भी हट जाए।
\end{solution}
